% Diskussion: 1-1 1/2 Seiten (Folie 16). Dem Dozenten besonders wichtig: Was bedeuten die Ergebnisse?
% Warum ist eine Methode besser? Auf welchen Daten funktioniert eine Variante?
% Alle Zahlen aus results/ms1/REPORT.md und results/why/REPORT.md. Regel: Ein Abstand zählt nur, wenn er
% größer ist als die doppelte größere Seed-Standardabweichung der beiden Methoden.
\section{Diskussion}
\label{sec:discussion}

\paragraph{Welche Methode ist die beste, und warum?}
LoRA $r=8$ ist nominell die beste Methode (96,21\,\%), doch Linear Probing (96,03\,\%) und LoRA $r=4$
(96,11\,\%) sind innerhalb der Seed-Streuung nicht von ihr zu unterscheiden. Nur volles Fine-Tuning
($-0{,}82$\,pp) und die CNNs ($-1{,}42$ und $-1{,}99$\,pp) sind klar schlechter. Die Zusatzexperimente
(\cref{sec:why}) legen nahe, warum ein eingefrorener Backbone genügt: Die vortrainierten Merkmale
trennen die Rassen bereits. Eine logistische Regression auf eingefrorenen ViT-B/16-Merkmalen erreicht
95,9\,\%, mit nur einem Bild pro Rasse schon 84,7\,\%. Zudem sind 22 unserer 25 Rassen
ImageNet-1k-Klassen~\cite{russakovsky2015imagenet}; bei Auswahl unter diesen 22 Rassen erreichen die
ursprünglichen ImageNet-Köpfe 96,2--96,8\,\%, ohne ein einziges Bild unseres Trainingssatzes gesehen zu
haben. Die Anpassung muss also keine Rassenmerkmale lernen, sondern vor allem vorhandene Merkmale auf
unsere Klassen abbilden, und das leistet ein linearer Kopf. Volles Fine-Tuning aktualisiert dagegen 86\,M
Parameter mit 2\,242 Bildern und landet leicht unter den effizienten Methoden. Das passt zur Beobachtung,
dass volles Fine-Tuning gute vortrainierte Merkmale verzerren kann~\cite{kumar2022finetuning}, auch wenn
Kumar \etal\ den Schaden vor allem außerhalb der Trainingsverteilung feststellen. Bei ResNet-50 zeigt sich
ein ähnliches Muster: Eingefrorene Merkmale mit Linear Probe erreichen 96,1\,\%, das vollständig
feinjustierte Netz 94,2\,\% (ein Lauf der Probe gegenüber drei Fine-Tuning-Seeds und ein anderes
Trainingsverfahren, der Vergleich ist daher nur ein Hinweis).

\paragraph{Lohnt sich LoRA? (A1, A2)}
A1 ist bestätigt: Mit 0,19--0,70\,\% trainierbaren Parametern ist LoRA nie schlechter als volles
Fine-Tuning, und LoRA $r=8$ ist sogar signifikant besser ($+0{,}82$\,pp bei einer Schwelle von
0,78\,pp); die anderen Ränge liegen innerhalb der Streuung. A2 ist widerlegt: Linear Probing liegt nur
0,17\,pp hinter der besten LoRA-Variante, innerhalb des Rauschens, trainiert dabei 16-mal weniger
Parameter als LoRA $r=8$ und ist am schnellsten fertig (im Mittel 0,3\,min; seine beste Epoche,
$3{,}7\pm1{,}2$, ist nominell auch die früheste). Auf diesem Datensatz gilt die Empfehlung der Vorlesung,
bei weniger als 10\,k Bildern nur den Kopf zu trainieren, sogar für feinkörnige Rassen. Gegenüber
Adapter-Methoden, die zur Inferenz zusätzliche Module brauchen, hat LoRA den Vorteil, dass die Änderung in
die Gewichte eingerechnet wird.

\paragraph{Rang und angepasste Layer (A3).}
Der erste Teil von A3 ist bestätigt, der zweite nicht. Die Ränge 4, 8 und 16 unterscheiden sich um etwa
0,25\,pp, und die Anpassung aller linearen Layer mit $r=8$ (1,35\,M Parameter) ist nicht besser als nur
Query/Value ($-0{,}39$\,pp, innerhalb des Rauschens). In Analogie zum niedrigen intrinsischen Rang der
nötigen Änderung, den Hu \etal~\cite{hu2022lora} für Sprachmodelle berichten, scheint zusätzliche
Kapazität im Adapter ungenutzt zu bleiben, wenn die eingefrorenen Merkmale bereits gut sind.

\paragraph{ViT gegen CNN (A4).}
Das beste CNN (ConvNeXt-T, 94,78\,\%) liegt 1,42\,pp unter dem besten ViT. Der Abstand ist kleiner als
2\,pp, übersteigt aber die Seed-Streuung (Schwelle 0,92\,pp); A4 ist daher im strengen Sinn widerlegt:
Die Accuracies liegen nah beieinander, unterscheiden sich aber systematisch. Unsere Probes auf
eingefrorenen Merkmalen deuten auf die Ursache. Mit Linear Probe sind die Merkmale von ResNet-50
(96,1\,\%, nur ImageNet-1k) so gut wie die von ViT-B/16 (95,9\,\%, ImageNet-21k und 1k). Der Abstand
scheint also von der Anpassung der CNNs durch volles Fine-Tuning zu kommen, das auch beim ViT die
schwächste Methode ist, und nicht von einem besseren Pretraining des ViT. Eine methodische Lehre betrifft
die Lernrate: Mit der ViT-Lernrate ($5\cdot10^{-5}$) war ResNet-50 untertrainiert (Validation-Accuracy
91,5\,\%, Validation-Loss 1,35); erst eine pro Architektur auf dem Validierungssplit gewählte Lernrate
macht den Vergleich fair (\cref{tab:lr}).

\paragraph{Auf welchen Daten würde eine andere Methode gewinnen?}
Unser Ergebnis hängt davon ab, wie nah die Zielaufgabe an den Pretraining-Daten liegt. Die drei Rassen
ohne ImageNet-Klasse erreichen mit Linear Probing einen etwas kleineren F1-Wert (93,8 gegenüber
96,2\,\%), vor allem wegen des American Bulldog, der mit ähnlichen Rassen verwechselt wird. Für Daten fern
von ImageNet (andere Bilddomänen, Klassen ohne Entsprechung im Pretraining) oder für deutlich größere
Datensätze dürften eingefrorene Merkmale nicht mehr genügen, und LoRA oder volles Fine-Tuning sollten
vorn liegen, wie Kornblith \etal~\cite{kornblith2019transfer} für Autos und Flugzeuge beobachten und wie
es die Strategietabelle der Vorlesung nahelegt. Training ohne Pretraining ist bei dieser Datenmenge keine
Alternative: ViT-B/16 erreicht nur 11,2\,\% und ResNet-50 31,1\,\%; der Vorsprung des CNN geht
vermutlich auf die induktive Verzerrung der Faltungen zurück.

\paragraph{Augmentierung und Datenqualität.}
\emph{none}, \emph{light} und \emph{strong} unterscheiden sich um höchstens 0,3\,pp, innerhalb des
Rauschens. Bei so guten Merkmalen und einem eingefrorenen oder niedrigrangig angepassten Backbone gibt es
wenig Overfitting, das Augmentierung verhindern könnte, und starke Augmentierung kann Rassemerkmale wie
die Fellfarbe verändern. Die Daten sind balanciert und enthalten keine exakten Duplikate. 8,6\,\% der
Testbilder haben ein sehr ähnliches Trainingsbild (CLIP-Ähnlichkeit $\geq0{,}95$); ohne sie sinkt die
Accuracy um 0,34\,pp, die Ergebnisse sind durch Fast-Duplikate also höchstens leicht geschönt. Die
verbleibenden Fehler konzentrieren sich: Die vier Bully-Rassen machen 16\,\% der Testbilder aus,
verursachen aber 59\,\% aller Fehler. Diese Rassen sehen sich auf Fotos sehr ähnlich und unterscheiden
sich vor allem in Größe und Proportionen, die ein Bild ohne Maßstab kaum zeigt.

\paragraph{Limitationen.}
Wir verwenden einen Datensatz, dessen Klassen großteils in ImageNet vorkommen; der Vorteil von Linear
Probing muss sich daher nicht auf andere Domänen übertragen. Nur die CNN-Lernrate wurde optimiert, die
ViT-Methoden nutzen je eine feste Lernrate. Drei Seeds und ein Validierungssplit mit 250 Bildern (ein
Bild entspricht 0,4\,pp; die Modellauswahl für Station~2 war ein Gleichstand) begrenzen die Auflösung
kleiner Unterschiede. Die Trainingszeiten wurden auf einer einzigen GPU gemessen.

\paragraph{Fazit und Ausblick.}
Für einen kleinen feinkörnigen Datensatz nahe an ImageNet erreicht ein neuer linearer Kopf auf einem
eingefrorenen ViT-B/16 96\,\% mit 0,02\,\% trainierbaren Parametern und dem kürzesten Training; LoRA
bringt keinen messbaren Gewinn, und volles Fine-Tuning ist leicht schlechter ($-0{,}8$\,pp). Die auf dem
Validierungssplit gewählte Methode, LoRA $r=4$, dient in Station~2 als Teacher, den wir in deutlich
kleinere ViTs destillieren.
